\uuid{IVdd}
\titre{Une rétropropagation de bout en bout}
\niveau{L3}
\module{Analyse de données}
\chapitre{Réseaux de neurones}
\sousChapitre{Rétropropagation}
\theme{Gradient, ReLU, mise à jour des poids}
\auteur{Maxime Nguyen}
\datecreate{2026-09-25}
\organisation{AMSCC}
\difficulte{3}

\contenu{
\texte{On considère un réseau à une entrée, un neurone caché et un neurone de sortie :
\[
h(x)=\operatorname{ReLU}(w_1x+b_1),\qquad
F(x)=w_2h(x)+b_2,
\]
avec $\operatorname{ReLU}(u)=\max(0,u)$. On dispose de l'unique donnée d'apprentissage $x=1$, de cible $t=1$, et l'on utilise la perte
\[
E(w_1,b_1,w_2,b_2)=\frac12\bigl(F(1)-t\bigr)^2.
\]
Les paramètres initiaux sont $(w_1,b_1,w_2,b_2)=(1,0,2,0)$.}

\begin{enumerate}
\item \question{Calculer $h(1)$, $F(1)$ puis la valeur initiale de $E$.}
\reponse{La pré-activation cachée vaut $w_1x+b_1=1>0$. Ainsi, $h(1)=1$, $F(1)=2$ et $E=\frac12(2-1)^2=\frac12$.}

\item \question{Calculer d'abord $\partial E/\partial w_2$ et $\partial E/\partial b_2$, puis $\partial E/\partial w_1$ et $\partial E/\partial b_1$ aux paramètres initiaux. On rappelle que $\operatorname{ReLU}'(u)=1$ pour $u>0$.}
\indication{Commencer par $\partial E/\partial F=F(1)-t$, puis remonter du neurone de sortie au neurone caché.}
\reponse{Au point initial, $F(1)-t=1$. Comme $F=w_2h+b_2$, on obtient
\[
\frac{\partial E}{\partial w_2}=(F-t)h=1,\qquad
\frac{\partial E}{\partial b_2}=F-t=1.
\]
La pré-activation cachée étant positive, sa dérivée ReLU vaut $1$. Par la règle de la chaîne,
\[
\frac{\partial E}{\partial w_1}=(F-t)w_2x=2,\qquad
\frac{\partial E}{\partial b_1}=(F-t)w_2=2.
\]}

\item \question{Effectuer une mise à jour simultanée des quatre paramètres par descente de gradient avec un pas $\eta=0{,}1$.}
\indication{Soustraire $\eta$ fois la dérivée correspondante à chaque paramètre initial.}
\reponse{Le nouveau vecteur de paramètres est
\[
(w_1^{\mathrm{new}},b_1^{\mathrm{new}},w_2^{\mathrm{new}},b_2^{\mathrm{new}})
=(0{,}8,-0{,}2,1{,}9,-0{,}1).
\]}

\item \question{Avec ces nouveaux paramètres, recalculer $F(1)$ et $E$. La mise à jour a-t-elle amélioré la prédiction ?}
\reponse{La nouvelle pré-activation cachée vaut $0{,}8-0{,}2=0{,}6$, donc
\[
F(1)=1{,}9\times0{,}6-0{,}1=1{,}04,
\qquad E=\frac12(1{,}04-1)^2=0{,}0008.
\]
La perte est passée de $0{,}5$ à $0{,}0008$ : la prédiction s'est rapprochée de la cible.}
\end{enumerate}
}
